
\definecolor{gapgreen}{HTML}{2E5B45}
\definecolor{gaprule}{HTML}{5C8F6E}
\definecolor{gapfill}{HTML}{E7F1E9}
\definecolor{gapink}{HTML}{1C3527}

\newtcolorbox{researchgap}[1][Research Gap]{%
    enhanced,
    colback=white,
    colframe=white,
    frame hidden,
    boxrule=0pt,
    coltitle=white,
    colbacktitle=gapgreen,
    fonttitle=\small\bfseries,
    title={#1},
    left=0mm,
    right=0mm,
    top=7.5mm,
    bottom=0mm,
    before skip=1.4em,
    after skip=1.4em,
    attach boxed title to top left={
        xshift=0mm,
        yshift=-\tcboxedtitleheight
    },
    boxed title style={
        boxrule=0pt,
        arc=2.2mm,
        outer arc=2.2mm,
        left=3mm,
        right=3mm,
        top=1.2mm,
        bottom=1.2mm
    }
}

\newtcolorbox{gapblock}{%
    enhanced,
    colback=gapfill,
    colframe=gaprule,
    coltext=gapink,
    before upper={\setstretch{1.12}},
    boxrule=0pt,
    leftrule=1.8pt,
    arc=2.5mm,
    outer arc=2.5mm,
    left=3mm,
    right=2.5mm,
    top=2.5mm,
    bottom=2.5mm,
    before skip=0.3em,
    after skip=1em
}

% ------------------------------------------------------------

\begin{researchgap}
\begin{gapblock}
\small\ttfamily
Current studies evaluate child safety in areas like benchmarking, age inference, multi-turn safety, and readability separately. However, it is not apparent whether a model's response improves in safety, clarity, and accessibility when the same request is made with different age signals, or whether these improvements hold during an ongoing conversation.
\end{gapblock}
\end{researchgap}